\documentclass[aspectratio=169, 14pt]{beamer}
\usepackage{stampfly_slides}

\lesson{8}{システムモデリング}{System Modeling}

\begin{document}

% ------------------------------------------------------------------
\begin{frame}
  \titlepage
\end{frame}

% ------------------------------------------------------------------
\begin{frame}{今日のゴール / Today's Goal}
  \begin{block}{目標}
    プラント伝達関数を導出し、モデルに基づいて P ゲインを設計する
  \end{block}

  \vspace{1em}

  \begin{enumerate}
    \item モータ+機体のプラント伝達関数 $G_p(s)$ を理解
    \item P 制御の閉ループ特性（$\omega_n$, $\zeta$）を導出
    \item $\zeta$ から $K_p$ を設計 --- L05/L07 の経験をモデルで裏付け
  \end{enumerate}
\end{frame}

% ------------------------------------------------------------------
\begin{frame}{プラントモデル / Plant Model}
  \begin{center}
    \vspace{-2mm}
    \resizebox{0.85\textwidth}{!}{\documentclass[border=10pt]{standalone}
\usepackage{tikz}
\usepackage{amsmath}
\usetikzlibrary{arrows.meta, positioning, shapes.geometric, calc}

\begin{document}

\definecolor{blockcolor}{RGB}{0, 120, 200}
\definecolor{arrowcolor}{RGB}{80, 80, 80}
\definecolor{plantcolor}{RGB}{40, 160, 40}
\definecolor{sumcolor}{RGB}{120, 80, 150}
\begin{tikzpicture}[
    block/.style={
      rectangle, draw=#1, line width=1.5pt,
      fill=#1!8, rounded corners=3pt,
      minimum width=38mm, minimum height=16mm,
      font=\bfseries, text=#1, align=center
    },
    arrow/.style={-{Stealth[length=3mm]}, line width=1.2pt, arrowcolor},
    lbl/.style={font=\small, above, text=arrowcolor},
  ]

  % Input
  \node[font=\bfseries, arrowcolor] (input) at (-1.5, 0) {$u$};

  % Motor block (actuator + 1st-order lag)
  \node[block=blockcolor] (motor) at (3.5, 0)
    {Mixer + Motor\\[2pt]$\dfrac{K_m}{\tau_m s + 1}$};

  % Body block (rigid-body integrator)
  \node[block=plantcolor] (body) at (10, 0)
    {Body\\[2pt]$\dfrac{1}{I \cdot s}$};

  % Output
  \node[font=\bfseries, arrowcolor] (output) at (14.5, 0) {$\omega$};

  % Arrows
  \draw[arrow] (input) -- (motor);
  \draw[arrow] (motor) -- node[lbl] {$\tau$} (body);
  \draw[arrow] (body)  -- (output);

  % Labels
  \node[font=\scriptsize, blockcolor, below=2mm] at (motor.south)
    {actuator dynamics};
  \node[font=\scriptsize, plantcolor, below=2mm] at (body.south)
    {rigid-body rotation};
  \node[font=\scriptsize, arrowcolor, above=1mm] at (input.north)
    {control input};
  \node[font=\scriptsize, arrowcolor, above=1mm] at (output.north)
    {angular rate};

  % Collapsed transfer function
  \node[font=\normalsize, arrowcolor, anchor=north] at (6.75, -2.8) {%
    $G_p(s) = \dfrac{K}{s(\tau_m s + 1)},\qquad K = \dfrac{K_m}{I}$};

\end{tikzpicture}
\end{document}
}
  \end{center}

  \vspace{-0.5em}

  \begin{exampleblock}{2つのブロック}
    \begin{itemize}
      \item \textbf{Mixer + Motor}: duty 入力 $\to$ トルク
            --- 1次遅れ $\tau_m \approx 20$\,ms
      \item \textbf{Body}: トルク $\to$ 角速度
            --- 剛体回転の積分器 $1/(I \cdot s)$
    \end{itemize}
  \end{exampleblock}
\end{frame}

% ------------------------------------------------------------------
\begin{frame}{線形化とパラメータ / Linearization \& Parameters}
  \begin{block}{ホバー近傍の小信号仮定}
    推力変動が小さいとき、モータ応答は 1 次遅れで近似できる
  \end{block}

  \vspace{0.5em}

  \begin{table}
    \centering\small
    \begin{tabular}{lllll}
      \hline
      \textbf{パラメータ} & \textbf{記号} & \textbf{Roll} & \textbf{Pitch} & \textbf{Yaw} \\
      \hline
      慣性モーメント & $I$ & 9.16e-6 & 13.3e-6 & 20.4e-6\,kg\,$\cdot$\,m$^2$ \\
      モータ時定数 & $\tau_m$ & \multicolumn{3}{c}{0.02\,s（共通）} \\
      実効プラントゲイン & $K$ & $\sim$102 & $\sim$70 & $\sim$19\,rad/s$^2$ \\
      \hline
    \end{tabular}
  \end{table}

  \vspace{0.5em}

  \begin{exampleblock}{なぜ軸ごとに $K$ が違う？}
    $K = K_m / I$ --- 慣性が大きい軸は応答が遅い。\\
    Yaw は反トルクで駆動するため $K_m$ 自体も小さい。
  \end{exampleblock}
\end{frame}

% ------------------------------------------------------------------
\begin{frame}{P 制御の閉ループ / P Control Closed Loop}
  \begin{center}
    \vspace{-2mm}
    \resizebox{0.85\textwidth}{!}{\documentclass[border=10pt]{standalone}
\usepackage{tikz}
\usepackage{amsmath}
\usetikzlibrary{arrows.meta, positioning, shapes.geometric, calc}

\begin{document}

\definecolor{blockcolor}{RGB}{0, 120, 200}
\definecolor{arrowcolor}{RGB}{80, 80, 80}
\definecolor{plantcolor}{RGB}{40, 160, 40}
\definecolor{fbcolor}{RGB}{220, 50, 30}
\definecolor{sumcolor}{RGB}{40, 160, 40}
\begin{tikzpicture}[
    block/.style={
      rectangle, draw=#1, line width=1.5pt,
      fill=#1!8, rounded corners=3pt,
      minimum width=30mm, minimum height=14mm,
      font=\bfseries, text=#1, align=center
    },
    sum/.style={
      circle, draw=sumcolor, line width=1.5pt,
      fill=sumcolor!8, minimum size=10mm,
      font=\large\bfseries, text=sumcolor
    },
    arrow/.style={-{Stealth[length=3mm]}, line width=1.2pt, arrowcolor},
    lbl/.style={font=\small, above, text=arrowcolor},
  ]

  % Input
  \node[font=\bfseries, arrowcolor] (input) at (-1.5, 0) {$r(t)$};

  % Sum
  \node[sum] (sum) at (1.5, 0) {$\Sigma$};

  % Controller
  \node[block=blockcolor] (kp) at (5.5, 0) {Controller\\$K_p$};

  % Plant
  \node[block=plantcolor] (plant) at (10, 0)
    {Plant\\$G_p(s)$};

  % Output
  \coordinate (out_end) at (14.5, 0);
  \node[font=\bfseries, arrowcolor] (output) at (14.5, 0) {};
  \node[font=\bfseries, arrowcolor, above=1mm] at (out_end) {$\omega(t)$};

  % Forward path
  \draw[arrow] (input) -- (sum);
  \draw[arrow] (sum) -- node[lbl] {$e(t)$} (kp);
  \draw[arrow] (kp) -- node[lbl] {$u$} (plant);
  \draw[arrow] (plant.east) -- ++(2.5, 0);

  % Feedback tap point
  \coordinate (fb_tap) at ($(plant.east)+(1.2, 0)$);
  \fill[arrowcolor] (fb_tap) circle (2.5pt);

  % Feedback path (T10 compliant: route below, connect to sum.south)
  \draw[arrow, fbcolor] (fb_tap) -- ++(0, -2.5)
    -| ($(sum.south) + (0, -0.5)$) -- (sum.south);

  % Sum signs
  \node[font=\scriptsize, sumcolor, anchor=south east]
    at ($(sum.north)+(150:0.4)$) {$+$};
  \node[font=\scriptsize, fbcolor, anchor=north east]
    at ($(sum.south)+(210:0.4)$) {$-$};

  % Closed-loop transfer function
  \node[font=\normalsize, arrowcolor, anchor=north] at (6.75, -3.8) {%
    $G_{cl}(s) = \dfrac{K_p K}{\tau_m s^2 + s + K_p K}$};

\end{tikzpicture}
\end{document}
}
  \end{center}

  \vspace{-0.5em}

  \begin{block}{閉ループ伝達関数}
    $G_{cl}(s) = \dfrac{K_p K}{\tau_m s^2 + s + K_p K}$
    \hfill $\longleftarrow$ \textbf{2 次系！}
  \end{block}
\end{frame}

% ------------------------------------------------------------------
\begin{frame}{2次系の特性 / Second-Order Characteristics}
  \begin{columns}[c]
    \begin{column}{0.5\textwidth}
      \begin{block}{固有振動数と減衰比}
        \begin{align*}
          \omega_n &= \sqrt{\frac{K_p K}{\tau_m}} \\[6pt]
          \zeta &= \frac{1}{2\sqrt{K_p K \tau_m}}
        \end{align*}
      \end{block}
    \end{column}
    \begin{column}{0.5\textwidth}
      \begin{block}{設計式: $\zeta$ から $K_p$ を逆算}
        \[
          K_p = \frac{1}{4\zeta^2 K \tau_m}
        \]
      \end{block}

      \vspace{0.5em}

      \begin{exampleblock}{L07 との接続}
        ステップ応答の\\オーバーシュート量 $\approx e^{-\pi\zeta/\sqrt{1-\zeta^2}}$
      \end{exampleblock}
    \end{column}
  \end{columns}
\end{frame}

% ------------------------------------------------------------------
\begin{frame}{ゲイン設計表 / Gain Design Table}
  \begin{table}
    \centering
    \begin{tabular}{lcccc}
      \hline
      & \textbf{$K$} & \textbf{$K_p=0.5$ の $\zeta$}
      & \textbf{$\zeta=0.7$ の $K_p$} & \textbf{$\zeta=1.0$ の $K_p$} \\
      \hline
      Roll  & 102 & 0.50 & 0.25 & 0.12 \\
      Pitch &  70 & 0.60 & 0.36 & 0.18 \\
      Yaw   &  19 & 1.15 & 1.34 & 0.66 \\
      \hline
    \end{tabular}
  \end{table}

  \vspace{0.5em}

  \begin{exampleblock}{読み取り}
    \begin{itemize}
      \item L05 の $K_p=0.5$ は Roll で $\zeta=0.50$（やや振動的）
      \item Yaw は $\zeta > 1$（過減衰 = 応答が遅い）--- 軸ごとに
            $K_p$ を変えるべき理由！
    \end{itemize}
  \end{exampleblock}
\end{frame}

% ------------------------------------------------------------------
\begin{frame}[fragile]{実習: モデルベースゲイン設計 / Hands-on}
  \begin{lstlisting}[style=sfcpp, title={user\_code.cpp (excerpt)}]
// Model-based Kp design (target zeta = 0.7)
const float tau_m = 0.02f;
const float K_roll  = 102.0f;
const float K_pitch =  70.0f;
const float K_yaw   =  19.0f;
float zeta = 0.7f;

float Kp_roll  = 1.0f/(4*zeta*zeta*K_roll *tau_m);
float Kp_pitch = 1.0f/(4*zeta*zeta*K_pitch*tau_m);
float Kp_yaw   = 1.0f/(4*zeta*zeta*K_yaw  *tau_m);

// Apply per-axis Kp in control loop
float roll_cmd  = Kp_roll  * (target_roll  - ws::gyro_x());
float pitch_cmd = Kp_pitch * (target_pitch - ws::gyro_y());
float yaw_cmd   = Kp_yaw   * (target_yaw   - ws::gyro_z());
  \end{lstlisting}
\end{frame}

% ------------------------------------------------------------------
\begin{frame}{チェックポイント / Checkpoint}
  \begin{alertblock}{確認事項}
    \begin{itemize}
      \item[$\square$] $G_p(s) = K / (s(\tau_m s + 1))$ を説明できる
      \item[$\square$] モデルベース $K_p$ と L05 の $K_p=0.5$ を比較飛行
      \item[$\square$] Roll/Pitch/Yaw の $\zeta$ 差を体感で理解
    \end{itemize}
  \end{alertblock}

  \vspace{1em}

  \begin{block}{次のレッスン}
    Lesson 9: 姿勢推定 --- センサフュージョンで角度を推定
  \end{block}
\end{frame}

\end{document}
